\chapter{قوانین انجمنی}
\label{ch:association}

\section*{اهداف این فصل}
\begin{itemize}
    \item درک مفهوم قوانین انجمنی و کاربرد آن در تحلیل سبد خرید
    \item تسلط بر معیارهای Support، Confidence و Lift
    \item یادگیری الگوریتم Apriori برای یافتن مجموعه‌های پرتکرار
    \item آشنایی با FP-Growth به‌عنوان روشی کارآمدتر
\end{itemize}

\section{مقدمه: از یادگیری با نظارت به بدون نظارت}

در چهار فصل گذشته، با وظایف \textbf{یادگیری با نظارت} (طبقه‌بندی و
رگرسیون) آشنا شدیم که در آن‌ها داده برچسب‌دار بود. از این فصل به بعد،
وارد دنیای \dmterm{\textbf{یادگیری بدون نظارت}}{Unsupervised Learning}
می‌شویم، جایی که هیچ برچسب از پیش‌تعیین‌شده‌ای وجود ندارد و هدف، کشف
ساختار یا الگوهای پنهان در خود داده است.

\dmterm{\textbf{قوانین انجمنی}}{Association Rules} یکی از قدیمی‌ترین و
شناخته‌شده‌ترین وظایف یادگیری بدون نظارت است. مثال کلاسیک آن
\dmterm{\textbf{تحلیل سبد خرید}}{Market Basket Analysis} است: با بررسی
میلیون‌ها رسید خرید یک فروشگاه، آیا می‌توان الگوهایی مانند «مشتریانی که
نان و کره می‌خرند، با احتمال بالا شیر هم می‌خرند» را به‌صورت خودکار کشف
کرد؟ چنین دانشی می‌تواند در چیدمان قفسه‌ها، طراحی تخفیف‌های ترکیبی و
سیستم‌های پیشنهاددهنده استفاده شود.

\section{مفاهیم پایه}

فرض کنید $I = \{i_1, i_2, \ldots, i_m\}$ مجموعهٔ همهٔ کالاهای موجود در
یک فروشگاه باشد و هر \dmterm{تراکنش}{Transaction} $T$ زیرمجموعه‌ای از
$I$ باشد (یعنی سبد خرید یک مشتری). یک \dmterm{\textbf{قانون انجمنی}}{Association
Rule} به‌صورت $A \Rightarrow B$ نوشته می‌شود، که در آن $A$ و $B$ دو
مجموعهٔ کالای مجزا از هم هستند (مثلاً $A = \{$نان، کره$\}$ و
$B = \{$شیر$\}$).

سه معیار زیر برای سنجش کیفیت یک قانون انجمنی به کار می‌روند:

\paragraph{پشتیبانی}
\dmterm{\textbf{پشتیبانی}}{Support} نشان می‌دهد چه نسبتی از کل تراکنش‌ها
شامل هر دو مجموعهٔ $A$ و $B$ با هم هستند:
\[
\text{Support}(A \Rightarrow B) = \frac{|\{T : A \cup B \subseteq T\}|}{|\text{کل تراکنش‌ها}|}
\]
پشتیبانی نشان‌دهندهٔ \textbf{شیوع} یک الگو در کل داده است؛ قاعده‌ای با
پشتیبانی بسیار کم، حتی اگر جالب باشد، ممکن است از نظر تجاری اهمیت کمی
داشته باشد چون به‌ندرت رخ می‌دهد.

\paragraph{اطمینان}
\dmterm{\textbf{اطمینان}}{Confidence} نشان می‌دهد از میان تراکنش‌هایی که
$A$ در آن‌ها وجود دارد، چه نسبتی $B$ را نیز شامل می‌شوند:
\[
\text{Confidence}(A \Rightarrow B) = \frac{\text{Support}(A \cup B)}{\text{Support}(A)}
\]
اطمینان معادل احتمال شرطی $P(B \mid A)$ است.

\paragraph{لیفت}
مشکل اطمینان این است که اگر $B$ در بین همهٔ تراکنش‌ها بسیار رایج باشد
(مثلاً اکثر مشتریان به هر حال نان می‌خرند)، اطمینان بالا می‌تواند
گمراه‌کننده باشد و نشان‌دهندهٔ یک رابطهٔ واقعی نباشد. برای رفع این مشکل،
از معیار \dmterm{\textbf{لیفت}}{Lift} استفاده می‌شود:
\[
\text{Lift}(A \Rightarrow B) = \frac{\text{Confidence}(A \Rightarrow B)}{\text{Support}(B)}
\]
اگر $\text{Lift} > 1$، حضور $A$ احتمال حضور $B$ را افزایش می‌دهد (رابطهٔ
مثبت). اگر $\text{Lift} \approx 1$، دو رویداد تقریباً مستقل‌اند (اطمینان
بالا صرفاً به دلیل رواج بالای $B$ بوده، نه یک رابطهٔ واقعی). اگر
$\text{Lift} < 1$، حضور $A$ احتمال حضور $B$ را کاهش می‌دهد (رابطهٔ
منفی).

\section{الگوریتم Apriori}

چالش اصلی در یافتن قوانین انجمنی، تعداد نمایی زیرمجموعه‌های ممکن از
کالاهاست: برای $n$ کالا، $2^n$ زیرمجموعهٔ ممکن وجود دارد که بررسی همهٔ
آن‌ها برای مجموعه داده‌های واقعی (با هزاران کالا) از نظر محاسباتی
غیرممکن است. \dmterm{\textbf{الگوریتم} \LR{Apriori}}{Apriori Algorithm}
با استفاده از یک ویژگی هوشمندانه به نام
\dmterm{\textbf{اصل Apriori}}{Apriori Principle} این فضای جستجو را
به‌شدت کاهش می‌دهد:

\begin{quote}
اگر یک مجموعه کالا پرتکرار\footnote{\LR{Frequent}} نباشد، هیچ‌یک از
ابرمجموعه‌های آن نیز نمی‌توانند پرتکرار باشند.
\end{quote}

به بیان دیگر، اگر مجموعهٔ $\{$نان، کره$\}$ پشتیبانی کافی نداشته باشد،
نیازی به بررسی $\{$نان، کره، شیر$\}$ نیست، چون پشتیبانی آن هرگز نمی‌تواند
از پشتیبانی زیرمجموعه‌اش بیشتر باشد.

\subsection{تولید مجموعه‌های پرتکرار}

الگوریتم Apriori به‌صورت تکرارشونده و لایه‌به‌لایه عمل می‌کند:

\begin{enumerate}
    \item ابتدا پشتیبانی همهٔ مجموعه‌های تک‌عضوی محاسبه می‌شود و آن‌هایی
          که از یک آستانهٔ حداقلی\footnote{\LR{Minimum Support}} پایین‌ترند،
          حذف می‌شوند.
    \item از ترکیب مجموعه‌های پرتکرار باقی‌مانده، مجموعه‌های دوعضوی
          کاندیدا ساخته می‌شوند و دوباره پشتیبانی آن‌ها بررسی می‌شود.
    \item این فرآیند برای مجموعه‌های سه‌عضوی، چهارعضوی و به همین ترتیب
          ادامه می‌یابد تا مجموعهٔ کاندیدای جدیدی باقی نماند.
\end{enumerate}

\subsection{استخراج قوانین}

پس از یافتن همهٔ مجموعه‌های پرتکرار، برای هر مجموعه، تمام تفکیک‌های
ممکن آن به دو بخش $A$ و $B$ بررسی می‌شود و اطمینان هر قانون محاسبه
می‌گردد. قوانینی که اطمینان آن‌ها از یک آستانهٔ حداقلی
\footnote{\LR{Minimum Confidence}} بیشتر باشد، به‌عنوان قوانین نهایی
گزارش می‌شوند.

\paragraph{ضعف اصلی Apriori}
ضعف عملی این الگوریتم، نیاز به بارها اسکن کردن کل مجموعه داده (یک بار
به ازای هر سطح از اندازهٔ مجموعه کاندیدا) است که برای داده‌های بسیار
بزرگ، هزینهٔ زمانی و حافظه‌ای بالایی دارد.

\section{الگوریتم FP-Growth}

\dmterm{\textbf{الگوریتم} \LR{FP-Growth}}{Frequent Pattern Growth} برای
رفع ضعف اصلی Apriori طراحی شده است. ایدهٔ اصلی آن، فشرده‌سازی کل
مجموعه داده در یک ساختار درختی به نام
\dmterm{\textbf{درخت} \LR{FP}}{FP-Tree} است، به‌گونه‌ای که تنها با دو
بار اسکن کل داده (به‌جای اسکن‌های مکرر Apriori)، می‌توان همهٔ مجموعه‌های
پرتکرار را با پیمایش این درخت استخراج کرد. به همین دلیل، FP-Growth در
عمل معمولاً به‌مراتب سریع‌تر از Apriori است، هرچند پیاده‌سازی آن پیچیده‌تر
است. برای اهداف آموزشی این کتاب، درک اصل Apriori اولویت دارد؛ FP-Growth
را می‌توان نسخهٔ بهینه‌شدهٔ همان ایده در نظر گرفت.

\section{ارزیابی و انتخاب قوانین جالب}

الگوریتم‌های یافتن قوانین انجمنی معمولاً تعداد بسیار زیادی قانون تولید
می‌کنند که اغلب آن‌ها بی‌فایده یا بدیهی‌اند (مانند «مشتریانی که نان
می‌خرند، احتمالاً نان می‌خواهند»). برای انتخاب قوانین واقعاً
\dmterm{\textbf{جالب}}{Interesting}، معیار لیفت (که پیش‌تر دیدیم) نقش
کلیدی دارد؛ قوانینی با لیفت نزدیک به ۱ معمولاً کنار گذاشته می‌شوند، حتی
اگر پشتیبانی و اطمینان بالایی داشته باشند. علاوه بر این، قضاوت متخصص
حوزه (مثلاً یک کارشناس بازاریابی که می‌داند کدام رابطه از قبل شناخته‌شده
و کدام واقعاً جدید است) نقش تکمیلی مهمی در فیلتر نهایی قوانین ایفا
می‌کند.

\section*{پیاده‌سازی در پایتون}

پیاده‌سازی الگوریتم Apriori و استخراج قوانین با کتابخانهٔ
\texttt{mlxtend} (توابع \texttt{apriori} و \texttt{association\_rules})
در بخش «فصل ۶» پیوست الف (\ref{app:python}) آمده است.

\section*{خلاصهٔ فصل}

در این فصل، اولین وظیفهٔ یادگیری بدون نظارت این کتاب، یعنی قوانین
انجمنی، را بررسی کردیم. با سه معیار کلیدی — پشتیبانی، اطمینان و لیفت —
آشنا شدیم و دیدیم چگونه اصل Apriori فضای جستجوی نمایی مجموعه‌های کالا
را قابل مدیریت می‌کند. همچنین به‌طور مختصر FP-Growth را به‌عنوان
جایگزینی کارآمدتر معرفی کردیم. در فصل بعد، به دومین وظیفهٔ اصلی یادگیری
بدون نظارت، یعنی خوشه‌بندی، می‌پردازیم.

\section*{تمرین‌ها}

\begin{enumerate}
    \item از میان ۱۰۰۰ تراکنش یک فروشگاه، ۲۰۰ تراکنش شامل «نان» و ۵۰
          تراکنش شامل هم «نان» و هم «کره» هستند. پشتیبانی و اطمینان قانون
          $\{$نان$\} \Rightarrow \{$کره$\}$ را محاسبه کنید.
    \item اگر پشتیبانی $\{$کره$\}$ در کل داده ۱۵ درصد باشد، لیفت قانون
          تمرین قبل را محاسبه کنید و نتیجه را تفسیر کنید (آیا نان و کره
          رابطهٔ مثبت، منفی یا مستقل دارند؟).
    \item اصل Apriori را با کلمات خودتان توضیح دهید و بگویید چرا این اصل
          فضای جستجو را به‌شدت کاهش می‌دهد.
    \item چرا یک قانون می‌تواند اطمینان بسیار بالایی داشته باشد (مثلاً ۹۵
          درصد) اما از نظر تجاری بی‌ارزش باشد؟ با یک مثال توضیح دهید.
    \item مزیت اصلی FP-Growth نسبت به Apriori از نظر تعداد بار اسکن کردن
          داده چیست؟
\end{enumerate}
